\chapter{Método de avaliação}\label{cap:metodo}
\section{Desenho e unidade de análise}
A pesquisa compara cinco configurações de recuperação no mesmo conjunto de consultas e candidatas. O fator variado é a contribuição facial/global: somente face, somente global e três combinações com pesos fixos. A variável principal de resposta é mAP; precisão e revocação nas posições 5 e 10 complementam a interpretação. A comparação por consulta é pareada, pois cada configuração responde à mesma consulta sob o mesmo gabarito.

A unidade recuperada é a fotografia completa. A presença de uma identidade anotada define a relevância para a consulta principal. Não se avalia apenas se o rosto mais semelhante pertence à pessoa, tampouco se a fotografia retrata o mesmo evento. A separação dessas unidades evita confundir detecção de faces, classificação de identidade e utilidade de uma lista de imagens.

O desenho é descritivo e delimitado. Não há sorteio de uma população de álbuns, otimização validada de pesos, treinamento novo ou teste de hipótese com independência presumida. Todas as vinte consultas compartilham galeria e contexto de coleção; duas compartilham a mesma fotografia-fonte. Essa dependência é preservada na discussão, sem interpretar as consultas como vinte ensaios independentes da população geral.

\section{Coleções e papéis}
\begin{table}[htbp]\centering\small
\caption{Coleções na execução atual e função de cada avaliação}\label{tab:bases}
\begin{tabularx}{\textwidth}{lrrX}\toprule
Coleção & Fotos locais & Consultas & Papel\\\midrule
Gallagher & 589 & 20 & Comparação principal de recuperação por identidade anotada\\
LFW & 13.233 & 1.672 & Verificação facial local, condicionada à indexação\\
INRIA Holidays & 1.491 & 500 & Verificação global com AP adaptada\\
Álbum familiar & 171 usadas & --- & Aceitação funcional privada; sem benchmark de identidade\\\bottomrule
\end{tabularx}
\fonte{manifestos da execução de fechamento e relatório de aceitação de 28/09/2026; as contagens de LFW e Gallagher são anteriores às exclusões de detecção descritas no texto}
\end{table}

Gallagher é a coleção principal. Seus arquivos locais incluem anotações de identidade e coordenadas dos olhos, usadas para construir o protocolo. A fonte de aquisição registrada pelo projeto é a página \emph{Gallagher Collection Person Dataset}, mantida no endereço do responsável em Cornell. A página não respondeu à reconsulta documental; isso é uma limitação da verificação atual de seus termos, não autorização de redistribuição. A descrição de \citeonline{Oh2015} também identifica o recorte de anotações faciais da coleção. Não se presume que toda pessoa visível tenha rótulo.

LFW foi concebida para estudar comparação de faces sob variações de captura, com protocolo de pares descrito por \citeonline{Huang2007}. O projeto usa uma adaptação para recuperação, apresentada separadamente. Holidays contém grupos de imagens da mesma cena ou objeto; sua página oficial identifica 1.491 imagens e 500 consultas \cite{Holidays,Jegou2008}. Essas diferenças de relevância impedem interpretar os mAPs das três bases como uma classificação comum de dificuldade ou qualidade.

\section{Galeria, gabarito e consultas Gallagher}
As 589 imagens globais definem a galeria elegível. As 931 linhas de anotação mapeadas para essa galeria determinam a presença das identidades. Para cada consulta, o conjunto relevante é formado pelas fotografias que contêm a identidade anotada, excluída a fotografia-fonte. O total nas vinte consultas é de 884 relações de relevância. Uma fotografia pode ser relevante para mais de uma pessoa; portanto, esse total não é uma contagem de fotografias distintas.

O índice facial contém 1.303 rostos em 587 fotografias. Duas fotografias não tiveram face detectada, mas permanecem na galeria e no gabarito. A regra é essencial: o detector não define quais fotografias são verdadeiramente relevantes. Condicionar o gabarito ao sucesso do método avaliado retiraria falhas do universo e alteraria o significado das métricas.

O protocolo ordena identidades pela frequência de fotografias anotadas e seleciona vinte identidades elegíveis. Exige ao menos três fotografias relevantes após retirar a fonte e utiliza uma consulta por identidade. Dentro de cada identidade, as candidatas seguem a ordem do nome da imagem e do índice anotado da face. Essa seleção é determinística, mas não aleatória nem representativa de todas as pessoas da coleção.

O recorte de consulta é construído a partir das coordenadas dos olhos, com escala registrada igual a 6,0 e limites de recorte compatíveis com a imagem. O ponto médio dos olhos é convertido para as coordenadas do recorte. A face aceita precisa ter uma caixa detectada que contenha esse ponto. Havendo mais de uma caixa elegível, o desempate utiliza distância ao centro, menor área e coordenadas. Não existe substituição pelo maior rosto vizinho.

A preparação registrou três recortes descartados antes de completar as vinte consultas. Quando uma candidata não satisfaz a validação, o motivo é registrado e a preparação pode examinar a próxima candidata da mesma identidade na ordem definida. Isso é diferente de trocar a pessoa-alvo dentro do recorte. Na avaliação de uma consulta já congelada, ausência de face contendo o ponto anotado produz erro explícito. O procedimento seleciona consultas para as quais foi possível obter uma representação facial válida; essa condição limita a extrapolação para consultas em que a pessoa não é detectada.

As consultas finais envolvem vinte identidades e dezenove fotografias-fonte. O componente facial utiliza o recorte associado à pessoa anotada; o componente global utiliza o descritor da fotografia-fonte completa já presente no índice global. A fonte é retirada tanto do conjunto relevante quanto das candidatas. Cada ranking tem, assim, 588 fotografias. Reutilizar o descritor da fonte para formar a consulta global não significa deixá-la entre os resultados.

\section{Escores, ausência de face e ordenação}
Seja $F_i$ o conjunto de vetores faciais da fotografia $i$, $f_q$ o vetor da pessoa da consulta e $g_q,g_i$ os descritores globais. A contribuição facial é definida por
\begin{equation}\label{eq:face}
 s_f(q,i)=\begin{cases}
 \displaystyle\frac{1+\max_{f\in F_i}c(f_q,f)}{2},&F_i\ne\varnothing,\\[4pt]
 0,&F_i=\varnothing.
 \end{cases}
\end{equation}
O componente global é $s_g(q,i)=[1+c(g_q,g_i)]/2$. A fusão é
\begin{equation}\label{eq:fusao}
 s_{\alpha}(q,i)=\alpha s_f(q,i)+(1-\alpha)s_g(q,i),
 \quad\alpha\in\{1;0;0{,}9;0{,}7;0{,}5\}.
\end{equation}
Uma face ausente recebe contribuição zero após o tratamento de disponibilidade, e não cosseno zero remapeado para 0,5. Essa distinção garante que ausência de detecção não seja confundida com um escore facial mediano.

O protocolo pareado mantém todas as candidatas em cada configuração, incluindo as sem face. Ordena os resultados por escore decrescente e, nos empates, por identificador crescente da imagem. A soma usa os pesos registrados; não há seleção posterior de valores para maximizar o teste. Os resultados ``somente face'' deste protocolo são, portanto, a configuração facial dentro do universo comum, e não uma busca que elimina previamente todas as imagens sem detecção.

\section{Métricas e cálculo das diferenças}
Seja $r_q(k)$ igual a um quando a fotografia na posição $k$ é relevante para a consulta $q$, e zero caso contrário. Seja $R_q$ o conjunto relevante depois das exclusões. Para $K\in\{5,10\}$,
\begin{equation}
 P@K(q)=\frac{\sum_{k=1}^{K}r_q(k)}{K},\qquad
 R@K(q)=\frac{\sum_{k=1}^{K}r_q(k)}{|R_q|}.
\end{equation}
O denominador de precisão é sempre $K$. A AP considera as $N_q$ candidatas do ranking integral:
\begin{equation}\label{eq:ap}
 AP(q)=\frac{1}{|R_q|}\sum_{k=1}^{N_q}P@k(q)\,r_q(k),
 \qquad mAP=\frac{1}{|Q|}\sum_{q\in Q}AP(q).
\end{equation}
As definições correspondem à média das precisões em posições relevantes descrita por \citeonline{Manning2008}. O avaliador retorna zero para conjuntos relevantes vazios, mas a preparação Gallagher exige relevantes e não produz esse caso nas vinte consultas. Relevantes não recuperados não acrescentam termos ao numerador; o denominador continua sendo o total relevante.

Por exemplo, em um ranking didático com três imagens relevantes nas posições 1, 3 e 5, a AP é $(1+2/3+3/5)/3\approx0{,}7556$. O exemplo apenas esclarece a fórmula e não é uma consulta experimental. AP igual a 0,95 não significa 95\% de probabilidade de uma fotografia conter a identidade. A medida sintetiza a ordenação em relação a um gabarito.

A comparação pareada calcula $\Delta AP_m(q)=AP_m(q)-AP_{\mathrm{face}}(q)$. Valores acima de $10^{-12}$ são ganhos, abaixo de $-10^{-12}$ são perdas e os demais são empates. A tolerância é aplicada antes de arredondar a apresentação. Reportam-se média, mediana e contagens, sem teste de significância ou intervalo de confiança baseado em independência não demonstrada.

\section{Protocolos auxiliares}
Em LFW, 13.233 imagens foram examinadas, e 48 sem detecção não entraram no índice. Restaram 13.185 fotografias e 16.058 rostos. O protocolo escolhe uma consulta por identidade com pelo menos duas imagens indexadas, mediante chave determinística SHA-256 com semente 20260915, totalizando 1.672 consultas. A consulta usa o vetor da maior face da imagem indexada, enquanto a relevância deriva dos rótulos da coleção. Essa regra difere da vinculação por olhos anotados de Gallagher e é uma limitação adicional da verificação auxiliar. Após excluir a fonte, existem 13.184 candidatas por consulta. O mAP local não é a acurácia do protocolo oficial de pares.

Em Holidays, a primeira imagem de cada grupo é a consulta e as demais imagens do grupo são relevantes. Há 500 consultas e 1.490 candidatas por ranking após excluir a fonte. A AP local utiliza a Equação~\ref{eq:ap}; o avaliador oficial integra trapézios na curva precisão--revocação, diferença já registrada no fechamento. Com um único relevante na segunda posição, a fórmula local produz 0,5, enquanto a integração oficial produz 0,25. Os resultados não são diretamente equivalentes.

A cópia local Holidays contém três pares byte a byte idênticos em grupos distintos. A origem dessa duplicação não foi determinada. O avaliador exclui a fonte por identificador/caminho, não todas as cópias por hash, ao passo que a interface também emprega SHA-256. Os avaliadores auxiliares utilizam ordenação vetorizada em PyTorch, sem garantia de desempate idêntico em todo equipamento. Essas ressalvas acompanham seus números e impedem comparações externas sem harmonização do protocolo.

\section{Execução, proveniência e verificações documentais}
A execução científica utilizada é \metodo{fechamento\_20260928\_cuda}, realizada no commit \metodo{91a2a28}, com árvore limpa registrada nos manifestos. O ambiente usa Python 3.10.20, InsightFace 0.2.1, ONNX Runtime GPU 1.23.2, PyTorch 2.11.0+cu128 e torchvision 0.26.0+cu128. A verificação de ambiente registra NVIDIA GeForce RTX 3050 Ti Laptop GPU, com aproximadamente 4 GiB de memória. Esses dados caracterizam a execução, sem constituir recomendação de hardware ou medida de desempenho.

A indexação global usa lotes de 32 imagens e quatro trabalhadores; a configuração geral registra avaliação auxiliar em lotes de 128. O protocolo pareado possui sua configuração própria, registrada no manifesto de fusão. Os hashes dos três pesos e das entradas/saídas permitem distinguir os artefatos efetivos de arquivos homônimos.

O fechamento registra treze manifestos, conferência de 92 referências a arquivos e hashes de 15.265 fotografias-fonte indexadas. Também registra a reconstrução do gabarito, a verificação geométrica das consultas e a preservação das referências anteriores. Trata-se de evidência já registrada. Nesta etapa de escrita foram reconferidos o manifesto final, as métricas utilizadas e a preservação dos arquivos da execução; as médias das cem linhas Gallagher e os deltas foram recalculados sem inferência.

Os rankings integrais foram usados pelo avaliador, mas somente o Top-10 foi persistido para inspeção. Portanto, recalcular a média das APs não é recalcular independentemente cada AP a partir de um ranking integral salvo. Essa fronteira é mantida na rastreabilidade do Apêndice~\ref{ap:reproducao}. Os números de execuções anteriores não foram mesclados com os atuais.

\section{Dados, privacidade e uso responsável}
Fotografias e representações faciais exigem cuidados de tratamento e divulgação. A LGPD é uma referência normativa para essa avaliação \cite{Brasil2018}; este trabalho não emite parecer jurídico nem presume dispensa ética. Execução local reduz a necessidade de transmitir imagens, mas não resolve automaticamente finalidade, elegibilidade, controle de acesso, retenção ou direitos das pessoas.

Código, pesos e fotografias são recursos distintos. A política consultada de InsightFace diferencia a licença MIT do código das condições de pesquisa não comercial dos modelos pré-treinados \cite{InsightFace}. A página Holidays atribui direitos das imagens ao INRIA \cite{Holidays}. A indisponibilidade das páginas de Gallagher e LFW impede afirmar nova confirmação integral de seus termos. O relatório original LFW permaneceu acessível em página do autor, o que permite fundamentar a tarefa sem substituir a verificação das condições atuais da coleção.

O pacote documental contém métricas agregadas e proveniência derivada; não contém fotografias, recortes, pesos ou embeddings. Para a comparação por consulta, os pares são calculados localmente antes de retirar os identificadores da apresentação. O gráfico mostra a distribuição completa de diferenças sem ligação entre pontos e pessoas. Isso reduz exposição e não é alegação de anonimização do acervo. A rastreabilidade detalhada permanece local, e a futura disponibilização da monografia exige confirmação institucional própria.
